Originally posted at
\url{http://vinipsmaker.wordpress.com/2013/08/13/splines-extraction-on-kopf-lischinski-algorithm-part-2/}
on \formatdate{13}{8}{2013}.

\subsection{Once upon a time\ldots{} a polygon with holes appeared}

The polygon-union described in \hyperref[blogpart1]{part 1} works correctly and
there are no better answers than the correct answer, but these polygons were
meant to be B-Splines and what have been done till now is insufficient.

Consider the following image.

\begin{figure}[H]
  \centering
  \includegraphics[width=0.8\textwidth]{assets/polygon-holes.pdf}
  \caption{Normal polygon (left) and bordered polygon (right)}
\end{figure}

The left polygon is the polygon what you see before and after the polygon-union.
They are visually indistinguishable. The right polygon has an outline to allow
you understand its internal representation. You can understand how the
polygon-union generates a polygon with a ``hole'' seeing the right polygon.

If the polygon-union don't generate visually identifiable pixels, then there
shouldn't exist any problem, but when we convert the polygon to a B-Spline, some
areas of the image won't be filled. The polygons won't fit correctly, like show
in the image below.

\begin{figure}[H]
  \centering
  \includegraphics[width=0.8\textwidth]{assets/polygons-with-holes-and-splines.pdf}
  \caption{After B-Spline conversion}
\end{figure}

The solution is to create some holes. With the representation of polygons used
in libdepixelize, this task is very simple and the key operation to accomplish
this goal is to find the ``edges that the polygon share with itself''. The
following paragraphs explain this task.

\begin{figure}[H]
  \centering
  \includegraphics[width=0.8\textwidth]{assets/polygon-holes-conversion.pdf}
\end{figure}

The above image has 2 pairs of hidden points: pair <5, 6>, that is equal to pair
<11, 10>, and pair <12, 13>, that is equal to pair <18, 17>.


The algorithm outline is:

\begin{enumerate}
\item Iterate over the points and, for each point, try to find another point
  that is equal to it. For each point
  \begin{enumerate}
  \item Get the internal sequence. To gather the internal sequence, just iterate
    over more nodes to get the ``largest common edge''.
  \item Make sure the ``largest common edge'' has a size greater than 1 (e.g. is
    not only a common vertex, but a common edge). Use this sequence to construct
    the internal ``hole'' (or holes, if more than one) polygon and then remove
    it from the original polygon.
  \end{enumerate}
\end{enumerate}

Beware that \emph{use a sequence to construct the internal ``hole'' polygon} is
another algorithm (and ain't trivial). It goes like this (let's name it
\emph{fill\_holes}):

\begin{enumerate}
\item Create an empty sequence of nodes. Let's name it \emph{hole}.
\item Let's name the beginning of our input sequence \emph{region\_begin}.
\item Iterate over the nodes of sequence and for each node:
  \begin{enumerate}
  \item Find another node equal to this one and if there is such node:
    \begin{enumerate}
    \item Let's name our new finding \emph{aux}.
    \item Append a copy of the range from \emph{region\_begin} until the current
      node (do \textbf{not} confuse the current node with \emph{aux}) into
      \emph{hole}
    \item We can see a common vertex: Current node and \emph{aux}. Find the
      largest common edge and there will be references to two more nodes. Call
      \emph{fill\_holes} with this new range.
    \item Forget about the range between current node and \emph{aux}. We won't
      see these points in the next iteration and \emph{region\_begin} will refer
      to \emph{aux} from now on.
    \end{enumerate}
  \end{enumerate}
\item Append the range from \emph{region\_begin} until the end of input
  sequence into \emph{hole}.
\item \emph{hole} is ready. Do something with it (e.g. store it).
\end{enumerate}

The \emph{fill\_holes} procedure is complicated and you should check
libdepixelize's source code (\emph{HomogeneousSplines::\_fill\_holes}) to kill
any ambiguity or doubt.

\begin{itemize}
\item Remark \#1: libdepixelize do the steps of the outmost algorithm (the one
  removing nodes from the sequence) in backward order, to avoid too much moves
  of elements in the internal vector.
\item Remark \#2: The polygon is represented in clockwise order, but the holes
  will be represented in counter-clockwise order, but there is no reason to
  normalize. You can safely ignore this (de)normalization feature.
\item Remark \#3: The reason behind the special procedure to make sure the
  common sequence is an edge and not a vertex come from the fact that the code
  to fix the position of   T-junction nodes creates a pattern where the polygon
  shares a point with itself.
\end{itemize}

\subsection{And then, the polygon was meant to be a B-Spline}

It's a chapter about splines extraction and it's fair to end with this step.

\begin{figure}[H]
  \centering
  \includegraphics[width=0.8\textwidth]{assets/evolution.pdf}
  \caption{It's evolution, baby}
\end{figure}

The algorithm is very simple. The points you get through the polygon-union
algorithm in \hyperref[blogpart1]{part 1} will be the control points of
quadratic Bézier curves and the interpolating/ending points of each Bézier curve
will be the middle points between two adjacent control points. This is the
special case where all points are smooth. The image below can help you get the
idea.

\begin{figure}[H]
  \centering
  \includegraphics[width=0.8\textwidth]{assets/points-location.pdf}
  \caption{Locations of the B-Spline points}
\end{figure}

If you have non-smooth points, then all you need to do is interpret them as
interpolating points, instead of control points. You'll have three interpolating
points that are part of two straight lines. There is a bit more of theory to
understand why the generated B-Spline is smooth.
